\section{Problem statement and scope}
\label{sec:motivation}

In crystalline solids, the interaction between electrons and lattice vibrations
controls electrical resistance, conventional superconductivity, thermoelectric
transport, carrier relaxation, and device heating.  DFPT is the most widely used
first-principles tool for calculating phonons and electron--phonon matrix elements
because it avoids a separate large-supercell calculation for every displacement
and treats the self-consistent electronic response directly
\citep{Baroni2001,Gonze1995,Giustino2017}.

The motivating observation of Quantum Harness Issue 35 is that DFPT is often
quantitatively useful in a surprisingly broad range of materials
\citep{Issue35}.  This statement needs careful wording.  It does \emph{not} mean
that ordinary semilocal DFPT is uniformly accurate for every correlated material,
every phonon, and every observable.  Standard DFPT can fail when the reference
electronic structure is qualitatively wrong, when a phonon couples to an
unprotected internal degree of freedom, or when the electronic response is
strongly frequency dependent.  The real puzzle is narrower and more useful:

\begin{keypoint}[Core scientific question]
Why are corrections beyond static Kohn--Sham linear response small for some
electron--phonon perturbations, including some perturbations in correlated
materials, while they are large for other modes in the same material?
\end{keypoint}

This formulation rules out two incomplete answers.

\begin{enumerate}
  \item \textbf{``The quasiparticle residue cancels the vertex.''}  A numerical
  product close to one is evidence, not a physical explanation.  The Ward identity
  produces such a cancellation only in a specific conserved-charge limit, whose
  order of limits is not the generic adiabatic phonon limit.
  \item \textbf{``DFPT works because the phonon shifts every electronic state by
  the same amount.''}  Real perturbations do not act equally on all Bloch states.
  Even a local operator proportional to the identity inside a correlated orbital
  subspace produces band- and momentum-dependent matrix elements because Bloch
  states have different orbital weights and phases.
\end{enumerate}

The proposal instead treats the one-body derivative of the electronic Hamiltonian
as an operator.  It separates the single global identity
\texttt{global\_charge}, relative block-identity shifts
\texttt{site\_charge}, block-internal traceless structure
\texttt{internal}, and inter-block matrix elements \texttt{nonlocal}.
Derivatives of interaction parameters are two-body vertices in a different
operator space.  Conservation laws can constrain only the first one-body
direction, and only when the original unprojected perturbation is independently
verified as a strict uniform $\vq=0$ full-space common shift.

\subsection{Target deliverable}

The target is a computational framework that uses DFPT as a baseline and adds a
small number of physically interpretable, channel-resolved many-body corrections.
It should be:

\begin{itemize}
  \item \textbf{simpler} than a full many-body calculation of every
  $g_{mn\nu}(\vk,\vq)$;
  \item \textbf{potentially cheaper than a dense beyond-DFPT vertex grid} by
  focusing expensive calculations on risky modes and low-dimensional operator
  channels, without claiming a measured speedup over DFPT;
  \item \textbf{transparent} because every correction is attached to a named
  physical operator and a documented validity test;
  \item \textbf{falsifiable} on materials for which DFPT succeeds for one mode
  but fails for another; and
  \item \textbf{uncertainty-aware}, including an abstention path when a mode lies
  outside the validated regime.
\end{itemize}

The first validation object may be the uniform electron gas (UEG), but it is not
the final theory.  The UEG contains only a scalar density channel and therefore
provides a clean baseline for understanding what is special about that channel.
The proposed method must then survive multi-orbital and correlated-material tests.

\subsection{What is and is not being corrected}

DFPT computes several related quantities: force constants, phonon eigenvectors,
Born effective charges, dielectric response, and electron--phonon matrix elements.
The main target here is the last quantity: the physical quasiparticle scattering
vertex associated with a phonon.  Corrections to equilibrium structure,
anharmonicity, or the phonon spectrum may also matter, but they are separate axes
of error.  The proposed benchmark will keep the following gates distinct:

\begin{enumerate}
  \item the Kohn--Sham ground-state and phonon baseline;
  \item the electron--phonon matrix element $g$;
  \item the observable assembled from $g$, such as a linewidth, mobility, or
  pairing kernel.
\end{enumerate}

Passing one gate does not establish the next.
